% =====================================================================
% Preamble --- Supplementary Information
%
% The supplement does not use IEEEtran. IEEE requires supplementary material
% to be a separate file and places no page limit on it
% (https://www.computer.org/csdl/journal/ta/write-for-us/15060, "There is no
% page limit on supplemental files"), and a single-column layout is easier to
% read for material that is mostly tables, decision rules and long figure
% descriptions. The main paper's two-column measure would break the wider
% tables here and gains nothing.
%
% Figures and tables are numbered S1, S2, ... per IEEE convention. The preprint
% numbered them "Extended Data Figure N", which is a Nature-family device with
% no IEEE counterpart, and additionally suppressed the caption label, so a
% reader following a main-text pointer had nothing to match it against.
% =====================================================================
\documentclass[11pt,a4paper]{article}

\usepackage[utf8]{inputenc}
\usepackage[T1]{fontenc}
\usepackage[english]{babel}
\usepackage{lmodern}
\usepackage{microtype}

\usepackage[margin=2.5cm]{geometry}
\usepackage{setspace}
\onehalfspacing

\usepackage{amsmath,amssymb}
\usepackage{graphicx}
\usepackage{booktabs}
\usepackage{multirow}
\usepackage{longtable}
\usepackage{array}
\usepackage{caption}
% Figure descriptions here run long by design --- they are the material the
% main paper's caption budget could not carry. Set them single-spaced and one
% size down so a description that outruns its page continues on the next one
% rather than overprinting the page number.
\captionsetup{font={footnotesize,singlespacing},labelfont=bf,
              justification=justified,singlelinecheck=false,
              skip=6pt}
\usepackage{capt-of}

\usepackage[hidelinks]{hyperref}
\usepackage[numbers,sort&compress]{natbib}
\bibliographystyle{IEEEtran}

% \bstctlcite passes the @IEEEtranBSTCTL entry in references.bib through to
% IEEEtran.bst. IEEEtran.cls defines it and the main paper gets it for free;
% this document is article class, so it needs the definition IEEEtran's own
% documentation gives. Without it the supplement would keep the em-dash
% repeated-author form that the main paper turns off, and the two reference
% lists would disagree on style.
% \@bstctlcite takes an OPTIONAL argument naming the aux stream, so \bstctlcite
% has to supply the default via \@ifnextchar. Defining it as
% \def\bstctlcite{\@bsphack\@bstctlcite} instead gives
% "Use of \@bstctlcite doesn't match its definition" followed by "Too many }'s".
\makeatletter
\def\bstctlcite{\@ifnextchar[{\@bstctlcite}{\@bstctlcite[@auxout]}}
\def\@bstctlcite[#1]#2{\@bsphack
  \@for\@citeb:=#2\do{%
    \edef\@citeb{\expandafter\@firstofone\@citeb}%
    \if@filesw\immediate\write\csname #1\endcsname{\string\citation{\@citeb}}\fi}%
  \@esphack}
\makeatother

\newcommand{\thd}[2]{\begin{tabular}[b]{@{}c@{}}#1\\#2\end{tabular}}

\usepackage{upgreek}
\newcommand{\uV}{\ensuremath{\upmu\mathrm{V}}}

\newcommand{\emoAll}{Emo-Soundscapes (all 1{,}213)}
\newcommand{\emoOrig}{Emo-Soundscapes (600 originals)}

% One figure per page, image first, description flowing beneath it. A float is
% one unbreakable box, so an image plus a 500-word description silently clips
% when the two together exceed the page; \captionof lets the description flow.
\newcommand{\msfigure}[3]{%
  \clearpage
  {\centering\includegraphics[width=\textwidth]{#1}\par}%
  \captionof{figure}{#3}\label{#2}%
}

% Figure descriptions, generated from reports/FIGURE_LEGENDS.md.
% Loaded as macros: \caption is a moving argument and \input inside it is
% fragile.
% Generated from reports/FIGURE_LEGENDS.md by scripts/63_legends_to_tex.py
% Do not edit here; edit the Markdown and regenerate.

\newcommand{\legendfigone}{%
\textbf{Arousal models transfer between corpora of the same sound domain far better than across it, under every representation tested.} \textbf{a} Generalisation stages for six regression algorithms (RF, XGBoost, GBDT, SVR, KNN, ElasticNet) over 122 spectro-temporal descriptors. Each line follows one algorithm from within-corpus cross-validation, through same-domain transfer (music $\leftrightarrow$ music), to cross-domain transfer (music $\leftrightarrow$ environmental); grey points are the individual corpus pairs behind each mean. Within-corpus values are GroupKFold-5 estimates grouped by source recording; transfer values fit on the whole of one corpus and are evaluated on the whole of another, with nothing refitted. Note that `\texttt{grouped by source recording'' is a property of the group identifiers each corpus supplies, and on the environmental corpus those identifiers are one per clip: 613 of its 1,213 clips are mixtures of 30 source recordings, so GroupKFold there is arithmetically identical to KFold and near-duplicate mixtures span folds. Panels \textbf{a}--\textbf{f} use the corpus as distributed, for comparability with prior work; Table 1 of the main text gives the leak-free counterparts, on which the within-corpus reference falls from 0.744 to 0.646 and median cross-domain transfer rises from 0.043 to 0.129. \textbf{b} Cross-corpus transfer matrix, averaged over the six algorithms; the diagonal shows within-corpus cross-validation. Green indicates positive rank correlation, red negative. \textbf{c} Directional asymmetry. For each music corpus, transfer outward to the environmental corpus and inward from it, averaged over algorithms; the dashed line marks the same-domain mean (0.59). The inward direction is uniformly sealed ($\rho$ $\le$ 0.09 for all three corpora), whereas the outward direction is not and increases from PMEmo ($-$0.04) through DEAM (+0.13) to Soundtracks (+0.24). Film soundtracks contain substantial atmospheric material --- sustained tones, texture, sound design --- which is a plausible account of that ordering but not one measured here. The barrier is \textbf{not symmetric}, and we report the asymmetry without an explanation for it. \textbf{d} Duration control. Music clips truncated from 30 s to a 6 s centred window, matching the environmental corpus; filled circles are 30 s, open circles 6 s. Same-domain transfer survives truncation (0.67 $\rightarrow$ 0.63) while cross-domain transfer stays at zero (0.03 $\rightarrow$ 0.05), excluding clip length as the cause of the collapse. \textbf{e} Leave-one-category-out generalisation inside the environmental corpus (n = 1,213), holding out each of the six Schafer soundscape categories in turn. The shaded band marks the mean cross-domain level. Generalisation to unseen content categories holds (mean $\rho$ = 0.52; range 0.20 for }quiet\texttt{ to 0.70 for }human`, five of six categories above 0.33), so the collapse in \textbf{a}--\textbf{b} is specific to the domain border rather than to any content shift. \textbf{f} The same operation costed twice, under five representations. For each representation the corpus is swapped and nothing is refitted; the grey mark gives the performance lost when the new corpus is in the same domain, the red mark the loss when it is not, both as a fraction of that representation's own within-corpus performance, and the printed value is the gap between them in percentage points. Points are medians over the corpus pairs of \textbf{a}. For the three uncontaminated families the loss is 20--31\% within the domain and 84--95\% across it --- hand-specified 20.2/93.9\%, AST 22.0/84.2\%, MERT 29.1/94.9\%, Wav2Vec2 30.7/90.0\%, a gap of 59 to 74 percentage points (Table 1). The result is where the collapse stands, not how high it is: every representation transfers between corpora and none crosses between domains, which disposes of the reading in which the embeddings simply fail to move between datasets. AST is the family that most ought to cross, since AudioSet contains large numbers of environmental categories; it raises cross-domain $\rho$ from 0.047 to 0.125, a factor of 2.7 that is still a sixth of what the same representation reaches within a corpus. CLAP is drawn apart and greyed. Its pretraining set contains Freesound, the source of the environmental corpus, an overlap declared before any of these numbers existed. It is not a generally stronger representation --- 0.620 on music-to-music against AST's 0.619 and 0.614 for the hand-specified descriptors --- and its lower barrier (57.3\%) appears only in comparisons that necessarily involve the corpus its pretraining data overlap. With a single environmental corpus available, genuine crossing cannot be separated from familiarity with the target; the barrier claim rests on the three uncontaminated families, and CLAP is used neither as support nor as refutation. Corpora: DEAM (music, n = 1,233), PMEmo (music, n = 717), Soundtracks (music, n = 360) and Emo-Soundscapes (environmental, n = 1,213). Soundtracks (Eerola \& Vuoskoski, 2011) differs from the other music corpora in collection site, decade, rating protocol (nine-point Likert scales over eight emotion terms rather than continuous valence--arousal sliders) and material selection (a balanced design of twelve target emotions $\times$ 30 excerpts). Its inclusion therefore tests whether within-music transfer reflects shared acoustic structure or merely a shared annotation tradition; transfer into it from the other music corpora is 0.56, and from the environmental corpus 0.09. Panels \textbf{a}--\textbf{e} rest on 12 ordered corpus pairs, each evaluated with the six algorithms: over the 36 same-domain values transfer averages 0.592 (s.d. 0.130) and over the 36 cross-domain values 0.054 (s.d. 0.161). Target is continuous perceived arousal rescaled to $-$1\ldots{}1. All audio is loudness-normalised to $-$23 LUFS before feature extraction, which is an admission criterion for this figure rather than a cosmetic step (Fig. S1). Metric is Spearman $\rho$ throughout. Source data are provided as a Source Data file.}

\newcommand{\legendfigtwo}{%
\textbf{Target-side labels recover about half of either corpus-swap gap, and label-free alignment recovers a quarter of the cross-domain one.} \textbf{a} Fraction of the gap recovered by k labelled target-corpus examples, for both kinds of corpus swap. The gap runs from zero-shot transfer (k = 0, defined as 0\%) to a model trained on the target corpus itself (100\%), so the denominator is the target's own within-corpus reference and not same-domain transfer --- the question an engineer faces is how close labelling k examples gets to labelling all of them. Line and points, median over the six ordered corpus pairs of each kind; band, interquartile range across those pairs. Across domains: 19\%, 31\%, 44\%, 57\% and 71\% at k = 10, 25, 50, 100 and 200. Same domain: 9\%, 16\%, 34\%, 50\% and 77\%. The two curves differ in shape rather than in whether a price exists; the cross-domain curve rises faster at small budgets because its zero-shot starting point is much lower, and the same-domain curve overtakes it at 200. \textbf{b} Alignment without labels, with its sham control \textbf{split by validity}. Correlation alignment (CORAL) recovers 25\% of the cross-domain gap with no target labels. The sham aligns the source not to the target but to an unrelated third corpus, which by construction carries no target-specific information --- except that with one environmental corpus available, the third corpus lies inside the *target's* own domain on three of the six ordered pairs, so the sham is a valid negative control on only half of them. Separated, the valid sham recovers $-$11\% and the invalid sham +25\%: the entire apparent sham effect is the invalid half. Subspace alignment is worse than doing nothing ($-$5\%); the axis is not clipped at zero, because clipping would turn ``actively harmful'' into ``small''. \textbf{c} CORAL decomposed. Matching per-dimension variance alone recovers 5\%; adding correlation structure adds 20 points more, 80\% of the total. That second step is target-specific by construction, which is the claim the pooled sham was supposed to test and could not. Every fraction is computed within an ordered corpus pair and only then pooled. Pooling first --- averaging the correlations and taking one ratio against one median ceiling --- mixes pairs whose baselines differ by more than the effect being measured. The environmental corpus here excludes its 613 mixture clips, which are combinations of 30 source recordings that the supplied group identifiers do not group; on the corpus as distributed they leak across folds and inflate both the target's within-corpus reference and the apparent size of the gap. The estimator is ridge regression on the hand-specified descriptors. The 200-label point rests on five of the six cross-domain pairs and four of the six same-domain pairs: Soundtracks (n = 360) cannot give up 200 labelled examples and still be evaluated. Restricting every budget to the pairs present at all of them changes the cross-domain values by at most 2 points and the same-domain values by up to 5, and does not change the ordering. \textbf{What this figure corrects.} The preprint drew a single curve, filtered to the cross-domain pairs, and reported it in the abstract and discussion as the *same-domain* price while describing the cross-domain boundary as one on which no purchase could be found. This legend stated the correct condition; the prose did not. The preprint also reported that the return per doubling collapses after 100 labels, from a constant 14--16 points to 5. On the leak-free corpus the increments are 10, 14, 10 and 14 --- no collapse, and the largest increment is the last. That claim is withdrawn, and with it any statement about where the curve saturates: it had not saturated at the largest budget tested. Metric is Spearman $\rho$; all audio is loudness-normalised to $-$23 LUFS. Source data are provided as a Source Data file.}

\newcommand{\legendednine}{%
\textbf{Feature-selection routes recover different descriptor classes depending on what they test.} \textbf{a} UpSet-style recovery pattern. The matrix marks which of four selection routes recovered each descriptor; connected points indicate a descriptor recovered by more than one route. The top bars count routes per descriptor and the left bars give each route's set size, split by feature class. The four routes are: SHAP importance with per-feature permutation FDR (q = 0.10, 200 permutations); preservation of the sign of the cross-domain Cohen's d; native shape-function importance from an Explainable Boosting Machine; and stability across a Rashomon set of 10 near-optimal models sampled from 60 ($\rho$ $\ge$ $\rho$\_max $-$ 0.02). \textbf{b} Share of each route's recovered set that consists of temporal-variability descriptors, against the 50\% chance line. The route that tests within-domain statistical significance returns 2 of 12 (17\%), below chance; the three routes that probe generalisation or model multiplicity return 5 of 8 (63\%), 6 of 10 (60\%) and 4 of 6 (67\%). The four routes differ in statistical principle --- post-hoc attribution, effect size, additive shape functions, and model multiplicity --- not in hyper-parameters, so \textbf{b} is a dissociation and not a shortfall of tuning. The route that is standard in the field returns, at a rate below chance, precisely the spectral-level descriptors that Fig. 1 and Fig. S8 show do not survive the domain boundary. What the panel licenses is therefore a negative: no attribution method here supports a statement about what will transfer, which is why every boundary in this work is crossed and measured rather than predicted from the descriptors a model relies on inside its own corpus. The companion requirement --- that an apparent overlap between attribution sets be read against the overlap that shuffled labels produce --- is in Fig. 5a,b. Source data are provided as a Source Data file.}

\newcommand{\legendfigthree}{%
\textbf{A representation change recovers a specific deficit that more of the same descriptors does not reach.} \textbf{a} Cross-validated Spearman $\rho$ for eight affect axes of the Soundtracks corpus (n = 360), using 122 spectro-temporal descriptors alone (open circles) and with 39 tonal descriptors added (filled). Axes are ordered by gain; numbers give the gain. Blue marks the three axes for which musical mode is theoretically decisive (happiness, valence, sadness) --- a grouping fixed a priori, not read off the result. The criterion stated for this test was that the gain concentrate on happiness, at more than twice the mean gain of the other seven axes and above 0.05 in absolute terms. It does, at +0.132 against a mean of +0.045, a ratio of 2.9. We can produce no artifact predating the analysis for that criterion and do not present it as preregistered. The ordering of gains follows mode-dependence: happiness, valence (+0.083) and sadness (+0.077) lead, while energy (+0.012) and anger (+0.025), which depend on level, roughness and rhythm, gain least. \textbf{b} Tonal descriptors alone against spectro-temporal descriptors alone. For happiness, 39 tonal descriptors ($\rho$ = 0.580) outperform 122 spectro-temporal ones (0.474); the deficit was one of representation rather than of dimensionality. \textbf{c} Mechanism check. Rank of five tonal descriptors among all 161, from random forest importance, for the happiness and tension axes (log scale). The two families dissociate: major-key strength and mode score rank 1st and 2nd for happiness but 3rd and 20th for tension, whereas interval consonance and sensory roughness rank 4th and 7th for tension but 117th and 76th for happiness. Mode carries perceived happiness; consonance and roughness carry perceived tension. The ordering was selected by the model from 161 descriptors, not imposed. The preprint presented this as a falsification test of a two-kind framework --- an information-limited gap that a representation closes. That framing is withdrawn. The framework is withdrawn, so there is nothing to falsify; and the test's necessary first half, that additional data does *not* close the happiness gap, was never measured, since no sample-size curve exists for any Soundtracks axis. What the result supports is narrower and still useful: a targeted representation change can recover a specific deficit, and which representation helps depends on which axis is being predicted. Bootstrap over out-of-fold predictions, resampled by stimulus, puts the happiness gain at +0.132 (95\% CI [+0.059, +0.209]), its excess over the mean of the other seven axes at +0.087 ([+0.018, +0.159]) and the tonal-over-spectro- temporal margin on happiness at +0.106 ([+0.007, +0.208]); the last two lower bounds sit close to zero, and because the bootstrap resamples predictions without refitting the models the intervals are narrower than a full resampling would give. Six of six algorithms show a larger gain on happiness than the mean of the other axes. Tonality is not the mechanism of the domain barrier: adding tonal descriptors lowers music-to-environmental transfer by 0.021 (18 paired comparisons, p = 0.038 uncorrected), in the predicted direction but about 4\% of a gap that runs from 0.591 to 0.115 on the leak-free environmental corpus. Interval-consonance weights are fixed a priori from the standard ordering of sensory consonance and are not fitted, so they add no degrees of freedom. All audio is loudness-normalised to $-$23 LUFS; cross-validation is GroupKFold-5 and each cell reports the best of six algorithms. Source data are provided as a Source Data file.}

\newcommand{\legendfigfour}{%
\textbf{An acoustic edit moves model output, reliably per excerpt for only one of three interventions and not consistently across model and material.} \textbf{a} Dose--response curves for three interventions applied to 80 source excerpts (40 per domain) at four dose levels, with dose 0 the unmodified control. Curves show the mean predicted probability of the low-arousal class under a model trained on the music corpus and a model trained on the environmental corpus; dotted lines mark the control level. The modelled quantity throughout is P(low arousal), so a rising curve means the edit made the excerpt read as \textbf{calmer}. Slowing spectral change raises the low-arousal probability in both domains; injecting transients lowers it in both. Envelope compression is retained as a failed manipulation. The curve is an average over 80 excerpts and cannot by itself distinguish an effect from a mean shift; \textbf{d} does that. \textbf{b} Manipulation check, run before any dose--response was read. Within-clip Spearman $\rho$ between dose and the targeted descriptor; point size gives the fraction of clips moving in the same direction as the dose, and the shaded band marks |$\rho$| < 0.5. A target descriptor counts as moved only if |$\rho$| $\ge$ 0.5 and at least 75\% of clips move in the same direction, a criterion fixed in advance. Transient injection meets it on both of its target descriptors and spectral smoothing on two of three; envelope compression meets it on one of three and is carried through the figure as a failed manipulation rather than as a null effect. \textbf{c} Specificity diagnostic. The five descriptors moved most by each intervention at maximum dose, in units of the training-set standard deviation. \textbf{d} Significance against reliability, per intervention and per domain model. Abscissa, the fraction of the 80 excerpts whose own slope across the four doses has the intended sign; ordinate, $-$log10 of a sign-flip permutation p on those 80 slopes. \textbf{The 50\% line is the statistic's arithmetic floor, not a chance level:} sign consistency is computed against the sign of the observed median, so it cannot fall below 50\% by construction, and under the null its median is 0.54. A point near 50\% therefore describes an average over excerpts that disagree, not an effect the intervention has on audio in general, however small its p value, and values of 57--60\% are close to the floor rather than above chance. Filled markers meet the reliability criterion used here --- sign consistency $\ge$ 65\% and p\_perm < 0.01 --- and open markers do not; we can produce no artifact predating the analysis for that threshold and do not present it as preregistered. Spectral smoothing meets it in both domains (median slope +0.015 music, +0.074 environmental; sign consistency 70\% and 71\%; p\_perm = 1 $\times$ $10^{-4}$ in both). Transient injection reaches significance (p\_perm = 1 $\times$ $10^{-3}$ music, 8 $\times$ $10^{-4}$ environmental) at 50\% and 59\% sign consistency, at or barely above chance, and is therefore reported as a mean shift with no reliable per-excerpt effect and is not built upon. Envelope compression is uninterpretable here, having failed \textbf{b} (p\_perm = 0.50 and 0.98). **These statistics pool, for each model, its own 40 training-domain excerpts with 40 from the other domain, and the preprint read the pooled result as showing that the direction of an edit crosses the domain boundary even though the predictions do not --- "what fails is calibration rather than sign". Disaggregated, that does not hold.** The music model shows the effect on both kinds of material (+0.099 on music excerpts, +0.012 on environmental, both p < $10^{-3}$). The environmental model shows it strongly on *music* excerpts (+0.080, 100\% sign consistency) and not on environmental excerpts --- the material it was fitted on, and the material the claim is about --- where the median slope is $-$0.019 in the wrong direction at 57.5\% sign consistency and p = 0.073. Three of four model $\times$ material cells move as intended and the fourth is the one the claim needs most, so the claim is withdrawn. Independently of the data it was incompatible with the metric: every transfer number in this work is a Spearman rank correlation, which is exactly invariant under any strictly monotone recalibration, so a purely calibrational deficit cannot move $\rho$ at all. The perturbed models are binary random-forest classifiers of low- versus high-arousal, one per domain, and are not the regression models of Fig. 1. Panel \textbf{c} bounds what may be claimed: no intervention moves its target descriptor most. Transient injection moves \texttt{spec\_contrast6\_dmean} by 9.0 SD while moving \texttt{onset\_rate} by 2.0 SD, ranking the target twelfth. Because the co-moving descriptors are themselves temporal-variability measures, these experiments support a claim about the variability *dimension* and not about any individual descriptor, and they are statements about model behaviour: no listener response to these edits was measured. An earlier version of this analysis averaged the 80 excerpts within dose and correlated the four resulting points; at n = 4 the smallest attainable two-sided p for a Spearman correlation is 0.083, so that statistic could not have reached significance whatever the data showed, and the excerpt-level test in \textbf{d} replaces it. All interventions are re-normalised to $-$23 LUFS afterwards so that loudness cannot drive the effect. Source data are provided as a Source Data file.}

\newcommand{\legendfigfive}{%
\textbf{Reference distributions, positive controls, and the comparison they license.} \textbf{a} Null distribution of the number of descriptors shared by the top-20 SHAP sets of three models under 20 label permutations. The observed value (4) falls at the null mean (3.75), p = 0.381. Shuffled labels produce the same overlap as intact ones, so the statistic cannot separate signal from noise and no claim is built on it. \textbf{b} Cross-validated AUC under the same permutations, confirming that permutation destroyed the signal (null mean 0.496) while the intact labels reach 0.936. The failure in \textbf{a} is therefore a property of the statistic, not of the data. \textbf{c} Positive controls for two physiological pipelines addressing the same question. In BIRAFFE2 the habituation control is significant in the direction opposite to the established effect ($\rho$ = +0.113 across trials, p = 0.002; skin conductance responses are expected to decline), and post-stimulus windows do not exceed random windows (p = 0.267); applied per participant, none of 94 passed both controls and habituation ran in the expected direction for only 35\%, so no usable subset exists and that pipeline is withdrawn rather than reported. In ds002721 all three controls pass: relative occipital alpha exceeds frontal alpha in 31 of 31 participants (p = 1.9 $\times$ $10^{-9}$); blink-locked averages reach 164 \uV{} frontally, 3.98-fold the occipital amplitude and far above random windows (p = 9.3 $\times$ $10^{-10}$); and a sound-onset response is present (\textbf{d}). Dashed line, p = 0.05. \textbf{d} Grand-average response at fronto-central electrodes (F3, Fz, F4, C3, Cz, C4), time-locked to sound onset, n = 31. Dark band, s.e.m.; grey band, the 95th percentile of a subject-wise sign-flip null with max-statistic correction across the epoch; green shading, the interval exceeding it (372--588 ms). Peak |t| = 8.78 at 452 ms, p = 0.0001 (10,000 permutations). The pre-stimulus interval is not significant (p = 0.54), excluding a slow drift. No classical N1 is present; the stimuli have sharp onsets (median 25 ms to half-amplitude, Supplementary Information), so the absent fast component is attributed to playback-latency jitter, which limits millisecond-scale analysis but not trial-wise spectral measures. \textbf{e} Stimulus-level intraclass correlation for eight self-report scales and 156 EEG measures, computed on the same 31 listeners and with the same estimator after within-listener z-scoring. The two sets are \textbf{not} computed on identical trials, as the preprint stated: the EEG measures use 1,240 trials over 307 stimuli and the self-report scales 1,080--1,121 trials over 285--298 stimuli. Points are individual measures, vertical bars the median. The EEG set spans four families: regional band power (26), magnitude-squared coherence between region pairs (50), imaginary coherency (50, robust to volume conduction) and electrode-pair asymmetry (30); coherence and asymmetry are included because the original analysis of these recordings located its effect there, and restricting the set to band power would test the conclusion on a representation that excludes the effect in question. Coherence is the strongest EEG family (maximum intraclass correlation 0.092, 95\% CI [0.037, 0.142], split-half 0.51) and band power the weakest (0.040, split-half 0.17), against 0.221 [0.155, 0.281] for the best-agreeing self-report scale on the same trials --- a 2.4-fold gap between two intervals that do not overlap. The figure states one requirement in two settings: a statistic is uninterpretable without the distribution it would take under the null (\textbf{a}, \textbf{b}), and a negative result is uninterpretable until the instrument is shown to detect a known effect (\textbf{c}, \textbf{d}). Panel \textbf{e} is the comparison those controls license, and it is reported as a bound rather than as an absence. The EEG value is the maximum of 156 measures and is biased upward for that reason, so it carries a null built the same way: permuting stimulus labels, recomputing all 156 intraclass correlations and taking the largest, 400 times, places the observed value at p = 0.02, while Benjamini--Hochberg across the same 156 does not pass it (smallest q = 0.16). The self-report value is also the maximum of a selection, but over 8 candidates rather than 156, so the two point estimates carry unequal selection bias in the direction that favours the EEG side; the comparison is conservative with respect to the conclusion drawn from it, but the asymmetry should be stated. The two corrections disagree because the value sits between them --- at rank one the Benjamini--Hochberg threshold is the Bonferroni threshold --- and 0.092 is close to the 0.10 this design resolves with 80\% power, so the interval is consistent with values from near zero up to about a tenth of the response variance. What the data support is the comparison and not a zero: whatever stimulus-specific structure the EEG carries is less than half of what the listeners' own ratings carry on the same trials. This is a statement about the component shared across listeners, not about whether the brain responds to sound --- the onset response in \textbf{d} is large and highly reliable; what is weak is the differentiating component. An earlier version of panels \textbf{c}--\textbf{e} was built on an electrodermal analysis that was retracted when its pipeline failed the controls now shown in \textbf{c}, and an earlier stimulus-level value of 0.892 was a leak --- the stimulus identifier itself had entered the measure set --- and is withdrawn. Both are documented in the Supplementary Information. Source data are provided as a Source Data file.}

\newcommand{\legendfigsix}{%
\textbf{The individual-level question is a design problem rather than a null result, and its price depends on the setting.} \textbf{a} Detection measured in real physiology. In an independent corpus of 75 listeners with about 240 auditory events each, the evoked response to a tone was tested per person by a sign-flip permutation test on the difference between real cues and randomly placed anchors. Each point gives the proportion of listeners in whom the response was detected when only n of that listener's observations were used (40 resamples per point). Three channels are shown, each with the analytic power curve for its measured effect size overlaid as a pale band; empirical and analytic differ by 0.015 on average. The dotted line is the false-positive rate obtained by sign-randomising the same differences, which preserves the noise and removes the effect; it stays at nominal (0.03--0.07) across the whole range, so the rise in the curves is detection and not false positives. The shaded band marks what a laboratory session supplies. Within it, an effect plainly visible in the group average ($-$5.6 \uV{} at the N1, group p = 0.0002 in all three channels) is recovered in 53\% of individuals from EEG, 39\% from pupil diameter and 23\% from heart rate. \textbf{b} Probability of detecting a within-person association as a function of the observations available per person, for four true association strengths. Synthetic data were generated from the empirical feature covariance of the EEG corpus (Ledoit--Wolf shrinkage) with an association of known strength injected, then passed through the identical pipeline, criterion and per-person permutation test used for the real data (500 simulations per cell, 100 permutations each, n from 30 to 5,000). Within the laboratory-session band, power is 0.10 for r = 0.40, 0.072 for r = 0.30, 0.06 for r = 0.20 and 0.04 for r = 0.10. The r = 0.10 curve is not monotone below n = 250; the Monte-Carlo s.e. is 0.009 to 0.015 per cell and the curve is included for the floor it establishes, not for its shape. \textbf{c} Observations per person required for 80\% detection, measured and simulated on one axis, against three supply levels. Every requirement is resolved by this grid; the two weakest effects, which a coarser run could only bound at > 1,000, now stand at 1,364 (r = 0.20) and 4,888 (r = 0.10). Dotted lines give the supply from this corpus (30), comparable public corpora (60), and a single overnight recording sampled once per minute over eight hours (480). Heart rate, the signal a consumer device records, requires roughly four times as many observations as scalp EEG (233 against 57). \textbf{d} When per-person calibration begins to pay, in three settings. For each descriptor--outcome cell the between-person standard deviation of the acoustic slope was estimated by restricted maximum likelihood (filled circle, median over that setting's cells) and compared with a null built by parametric bootstrap from that cell's own design (cross, median over the same cells). The null is plotted rather than subtracted, because a between-person spread estimated from few observations per person is biased upward: at 18 observations per person the estimator returns 0.11 when the true value is zero, and a reader who cannot see that cannot judge the physiological row. Break-even, n*, is the number of observations per person at which the corrected spread exceeds the error in estimating it. Urban soundscape ratings (42 observations per person, 8 cells): estimate 0.134 against a null of 0.013, n* = 67, and 2 of 8 cells sit above break-even at the observation count the corpus supplies, both on the eventfulness axis. Two qualifications attach to those two cells and neither appeared in the preprint. The corpus's own manipulation check passes for the pleasantness axis (p = 4 $\times$ $10^{-6}$) and \textbf{not} for eventfulness (p = 0.38), so the axis carrying both crossings has no passing positive control; and no out-of-sample personalisation test exists for this corpus, so ``above break-even'' here is a property of a variance model rather than a demonstration that a per-person slope beats the population slope. Music ratings (44 per person, 10 cells): 0.074 against 0.000, n* = 166, none across. Music electrodermal responses (18 per person, 15 cells): 0.106 against a null of 0.114 --- the estimate does not exceed its own noise --- with n* = 318 taken over the 5 cells in which it is finite, 11 of the 15 indistinguishable from their own null, and none across. Both kinds of estimate in \textbf{a}--\textbf{c} are lower bounds, for different reasons. The simulation assumes multivariate normality, a linear association and within-person stationarity, each of which favours detection. The measurement uses a deliberately plain pipeline --- a fixed-window mean over one electrode cluster, no independent component analysis and no spatial filtering --- so a dedicated pipeline would do better; it is, however, representative of what a non-specialist pipeline achieves. The extrapolation behind \textbf{c} was validated out of sample: predicting detection at n from an effect size estimated on a disjoint half of the same person's trials matches the observed rate to within 0.03 on average. Panel \textbf{d} is a different quantity from \textbf{a}--\textbf{c} --- break-even for per-person calibration, not detection of a within-person association. The three settings in \textbf{d} differ in corpus, descriptor set and outcome axis simultaneously (ISO-532 psychoacoustic descriptors against ISO soundscape axes; spectro-temporal and tonal descriptors against valence and arousal; the same descriptors against electrodermal measures), so the spread between them cannot be attributed to any one of the three and they are not compared against one another. Where the direct out-of-sample test was run it points the other way from the model: fitting each listener their own slope loses to the population slope in 15 of 15 electrodermal cells ($-$0.109 to $-$0.072 at 18 observations per person) and in every music-rating cell at 18, 30 and 50 observations per person, turning positive in 4 of 10 cells only at n = 100. Two earlier claims from that analysis are withdrawn and do not appear here: a between-person spread of 0.106 for the physiological setting, which sat below its own null; and the reading of a confidence interval excluding zero as evidence of individual differences, since that interval is for the composite of true spread and null and, at 18 observations per person with a true spread of zero, excludes zero on 84\% of occasions. Read together, \textbf{a}--\textbf{c} place the shortfall at roughly a factor of twenty and locate it in duration and repetition rather than in sample size: the quantity in short supply is observations within a person, which recruiting more participants does not provide. Source data are provided as a Source Data file.}

\newcommand{\legendfigeight}{%
\textbf{Withdrawn.} The preprint's summary figure placed the four boundaries on two axes --- the fraction of an attainable ceiling that survives each crossing, and the price of buying the gap back --- with colour carrying a diagnosis: budget-limited where target-side observations close the gap, information-limited where they do not. It is withdrawn rather than corrected, for three independent reasons. The diagnosis it encoded is withdrawn: the boundary the figure coloured information-limited is the one whose budget curve the paper reports, and that curve recovers 57\% of the gap at 100 target labels. Panel \textbf{a} placed four incommensurable ratios on one axis labelled as a fraction of an attainable ceiling. The corpus row is transfer against a within-corpus model reference; the response-channel row is one intraclass correlation against another, which is a ratio of reliabilities and not of achieved performance; the individual row is observations supplied against observations required; and the synthesis row was drawn at 1 and 0 by construction and was never a measured fraction at all. Panel \textbf{a}'s electrodermal row additionally used the wrong corpus's observation count as its numerator. Nothing in the paper now needs a figure of this kind. The quantities it summarised are reported per boundary, in the units each was measured in, and the ledger of what improved, what did not improve under the routes tested, and what the evidence cannot decide is given as text in the Discussion.}

\newcommand{\legendedone}{%
\textbf{Loudness acts as a corpus fingerprint rather than as signal.} \textbf{a} Cross-corpus AUC (DEAM $\rightarrow$ PMEmo) for five classifiers before and after loudness normalisation to $-$23 LUFS. Normalisation raises transfer for every model, from a mean of 0.866 to 0.886. \textbf{b} Mean SHAP rank of the five level-dependent descriptors before and after normalisation; arrows run from raw to normalised. All five lose importance, with \texttt{rms\_p90} falling from rank 39 to rank 93. Removing absolute level improves rather than degrades cross-corpus transfer, which is the signature of a corpus-specific confound: mastering practice differs systematically between corpora. Because every transfer number in this work is computed on normalised audio, this control is the admission criterion for the domain barrier, and the obvious objection --- that normalisation creates the barrier by discarding a genuinely informative dimension --- was tested separately. Three feature pipelines that share no data, no intermediate table, no constant and no estimate --- normalised descriptors, descriptors computed with loudness retained, and the 6,373-dimensional ComParE set extracted from the original recordings by the corpus authors --- put the best cross-domain transfer at $\rho$ = 0.045, 0.049 and 0.072. The spread of 0.027 is below the 0.05 threshold fixed before the test, and loudness alone contributes only +0.059. This is the only result in this work for which agreement between separately derived pipelines is claimed; no other comparison here meets that condition. Source data are provided as a Source Data file.}

\newcommand{\legendedtwo}{%
\textbf{Regression framing recovers the discarded middle of the arousal range.} \textbf{a} Clips available, used under the regression framing, and used under the original binary framing, per corpus. \textbf{b} Totals across corpora. The binary framing selects the two extremes of the arousal axis and discards the middle, retaining 1,341 of 3,782 annotated clips. The regression framing uses 3,163. The discarded clips carry the threshold information that a rule library requires. Source data are provided as a Source Data file.}

\newcommand{\legendedthree}{%
\textbf{Glass-box models cost almost no accuracy.} \textbf{a} Spearman $\rho$ under GroupKFold-5 for black-box and glass-box models in both domains. \textbf{b} Change in $\rho$ from imposing monotone constraints on the seven cross-domain-invariant descriptors, relative to the same feature subset without constraints. Monotone constraints cost +0.0003 in the ambient domain and $-$0.0016 in the music domain --- replicated at zero cost in both. The monotonicity assumed by threshold rules is therefore supported by the data rather than imposed on it. The fully additive EBM trails the best black box by 0.019. Source data are provided as a Source Data file.}

\newcommand{\legendedfour}{%
\textbf{Explainable Boosting Machine shape functions give thresholds without post-hoc attribution.} Additive contribution to predicted arousal as a function of descriptor value for the six most important terms in the ambient-domain EBM (n = 1,213, interactions disabled). Red lines mark the zero crossing, printed above each panel; these are the thresholds a rule library reads directly, with no attribution step. Five of the six terms are temporal-variability measures. Source data are provided as a Source Data file.}

\newcommand{\legendedfive}{%
\textbf{Descriptor importance is stable across the Rashomon set.} Share of near-optimal models ranking each descriptor in their top 15. The Rashomon set comprises the 10 of 60 sampled models within 0.02 Spearman $\rho$ of the best ($\rho$ $\in$ [0.823, 0.843]); models were sampled over algorithm family, depth, regularisation and subsampling, and importance was measured by permutation importance on a held-out fold. Six descriptors are recovered by every near-optimal model, four of them temporal-variability measures. Reporting importance from a single fitted model would have left this multiplicity untested. Source data are provided as a Source Data file.}

\newcommand{\legendedsix}{%
\textbf{No EEG measure tracks self-reported emotion across listeners at the effect size this design can resolve.} Each point is one of 208 combinations of an EEG measure (26) and a rating scale (8). The abscissa gives the mean across the 31 listeners of the within-listener partial Spearman $\rho$ between that measure and that listener's own trial-by-trial rating, with trial index partialled out; both EEG and ratings can drift over a session, and a shared drift would otherwise manufacture a correlation. The ordinate gives the evidence against a zero group mean (Wilcoxon signed-rank). Point size is the fraction of listeners whose sign matches the group mean; colour is the rating scale. The dashed line marks the Benjamini--Hochberg threshold at FDR < 0.05. No combination crosses it (largest |mean $\rho$| = 0.115; median 0.028; the group mean's sign is shared by a median of 55\% of listeners), and the distribution is symmetric about zero as expected under the null. This is a bound and not an absence. Listeners contributed a median of 30 rated trials each, and the design simulation (Fig. 6b) shows that at 30 observations per person a within-person association of r = 0.40 would be detected on 10\% of occasions and one of r = 0.20 on 6\% --- that is, at chance. The panel therefore establishes only that no large within-person association is obvious, and supports no conclusion about associations below approximately r = 0.5. The group-level result in Fig. 5e does not share this limitation, because it pools 31 listeners at the level of the stimulus. This figure is computed on the pipeline validated in Fig. 5c,d, which is what makes the negative interpretable at all; independent-component removal of ocular artefacts, used here as a sensitivity manipulation check, increased detection of the known effect while further reducing the effect of interest. An earlier version used an electrodermal analysis that was retracted when its pipeline failed the same controls. Source data are provided as a Source Data file.}

\newcommand{\legendedseven}{%
\textbf{A pretrained representation gains most on the one corpus that overlaps its pretraining data.} Best cross-validated Spearman $\rho$ for 122 hand-crafted spectro-temporal descriptors and for 512-dimensional pretrained audio embeddings, on three corpora (best of six algorithms, GroupKFold-5, loudness-normalised audio). Embeddings win on all three: +0.040 on DEAM (0.736 $\rightarrow$ 0.776), +0.034 on Soundtracks (0.823 $\rightarrow$ 0.856) and +0.070 on Emo-Soundscapes (0.773 $\rightarrow$ 0.843). The advantage does not require a strong downstream model --- on embeddings, cross-validated linear ridge regression is competitive with the tuned tree ensembles --- so it is carried by the representation and not by the estimator. The two music corpora have no overlap with the embedding model's pretraining data and give the two smaller margins; the corpus that derives from the public sound repository appearing in that pretraining data gives the largest. We report the music-corpus figure, +0.037, as the estimate that generalises and leave the difference unattributed. It cannot be assigned to pretraining overlap here: the one corpus sharing data with pretraining is also the only environmental corpus in the comparison, and both clean corpora are music, so overlap and sound domain are collinear. The larger margin is equally consistent with the embedding model simply being stronger on environmental sound. Separating the two accounts would require an environmental corpus outside the pretraining data, or a music corpus inside it. This figure concerns within-corpus accuracy only. Whether a pretrained representation crosses the domain boundary is a separate question, answered for four pretraining paradigms in Fig. 1f and Table 1, where none of them does. Source data are provided as a Source Data file.}

\newcommand{\legendedeight}{%
\textbf{Only temporal-variability descriptors keep their direction across sound domains.} \textbf{a} Standardised effect size (Cohen's d) of each FDR-significant descriptor in the music domain versus the environmental domain, with marginal densities by feature class. Shaded quadrants mark sign preservation; filled circles keep their sign across domains, open circles reverse it. The dotted diagonal is the identity line. Descriptors with min|d| > 0.35 are labelled. \textbf{b} The eight sign-preserving descriptors, ranked by the weaker of their two effect sizes (printed at the right of each row). Circles give |d| in each domain and the connecting bar gives the gap between them. All eight are negative in both domains: lower descriptor value corresponds to lower arousal, whether the sound is music or an environmental recording. The descriptors that reverse are without exception measures of absolute spectral level --- the property on which the two domains differ most. Feature classes are assigned by statistic type --- descriptors ending in \texttt{\_std} or \texttt{\_dmean}, plus \texttt{onset\_rate}, are temporal-variability measures; all others are spectral-level measures. Significance was established per domain by per-feature permutation testing (200 permutations) with Benjamini--Hochberg control at q = 0.10. This analysis is the second of the four selection routes in Fig. S9, and is the only one that tests the boundary directly. Source data are provided as a Source Data file. This analysis appeared as a main figure in an earlier version; it was moved to the Supplementary Information when the price of the first boundary (Fig. 2) took the main-figure slot.}


\date{}
